\section{A two-scale density profile inside a simple clearance caustic}
\label{sec:v65-simple-caustic-profile}
We now compute the density inside a part of the retained caustic,
including a roof-critical point where no roof pivot is available.
The conclusion is local on a simple competing-hit seam. It is not a
uniform classification of the entire caustic, and it does not convert
a tube-mass estimate into height control.

\subsection{A simple physical-side critical contact}
Fix $m,R$, an original center word, and a physical-side limiting
critical point $z$ of its reduced action. Assume every selected
contact has positive incidence. On an endpoint neighborhood $U$,
assume the only active physical boundary is one competing disk for
flight $j$, with its perpendicular foot strictly inside that flight.
All other first-hit inequalities and all section decisions are strict.
Use the same responsible disk and the image mark $j+1$ as in
Section~\ref{sec:v49-physical-layers}. The local signed margin is
\begin{equation}\label{eq:v65-simple-margin}
 g(x)=|Z_d(T_R^{j+1}x)|-R,\qquad g(z)=0.
\end{equation}
After shrinking $U$, the original word occupies exactly $g>0$ in
this neighborhood and its exact half-word labels are constant there.
The equality side has zero source measure. The local continuation of
the selected contact equations merely supplies coordinates for this
one-sided source.

The derivative $dg(z)$ is nonzero. Indeed
$\partial_\varphi Z_d=L_d\ge\ell_*>0$ on the selected envelope,
by \eqref{eq:v49-transverse-derivatives}; the regular recentering and
endpoint maps are local diffeomorphisms. Define
\begin{equation}\label{eq:v65-profile-scale}
 H=\nabla^2F(z),\quad t_z=F(z),\quad
 a_z=dg(z)H^{-1/2},\quad \kappa_z=|a_z|>0,\quad
 J_z=\frac{2\pi w_z}{\sqrt{\det H}}.
\end{equation}
Here $w_z=|F_{uv}(z)|/(4\pi R c_*)$ is the density of the original
section source in endpoint arclengths. Thus $J_z$ is given by
Proposition~\ref{prop:v65-boundary-monodromy}. At a boundary point
it is the full local quadratic coefficient; the physical one-sided
coefficient will be one half of it.

\begin{lemma}[Uniform angular threshold comparison]
\label{lem:v65-angular-threshold}
For $a\in\R^2\setminus\{0\}$ put $\omega_\theta=(\cos\theta,\sin\theta)$.
Suppose $q_r$ on the circle satisfies
$\|q_r-a\cdot\omega\|_\infty\le Lr$, $0<r<r_0$.
Uniformly in $\lambda\ge0$,
\begin{equation}\label{eq:v65-angular-threshold}
 \left|\bigl|\{0<q_r<\lambda\}\bigr|
       -2\arcsin\!\left(\min\{1,\lambda/|a|\}\right)\right|
       \le C_a\sqrt{Lr},
\end{equation}
with an enlarged constant when $Lr$ is not small. Bars on sets mean
angular Lebesgue measure. No differentiability of $q_r$ is required.
\end{lemma}
\begin{proof}
Rotate $a$ to $(|a|,0)$. The angular measure of
$\{0<|a|\cos\theta<\lambda\}$ is the displayed arcsine expression.
The symmetric difference of the two threshold sets lies in
\[
 \{|a|\cos\theta|\le Lr\}\ \cup\
 \{\bigl||a|\cos\theta-\lambda\bigr|\le Lr\}.
\]
For every real $b$ and $0<\eta\le1$, the set
$\{|\cos\theta-b|\le\eta\}$ has measure at most $C\sqrt\eta$.
To see this, split into the arcs within distance $2\sqrt\eta$ of
$0$ and $\pi$, whose total length is $O(\sqrt\eta)$, and their
complement, where $|\sin\theta|\ge c\sqrt\eta$. On each monotone
complementary arc the inverse function theorem bounds the length by
$C\eta/\sqrt\eta$. If the specified level is outside $[-1,1]$,
the same bound or the empty set applies. Rescale by $|a|$ and bound
the two pieces of the symmetric difference. This also proves uniformity
at the transition value $\lambda=|a|$, where a linear error bound
would not follow from this argument.
\end{proof}

Let $b_{z,s}$ denote the roof density of the original local source
with the additional restriction $0<g<s$. Its neighborhood is chosen
below as a full small Morse disk, restricted to the physical side.
This density is not renormalized.

\begin{theorem}[The simple seam--minimum profile]
\label{thm:v65-seam-profile}
There are $h_z>0$, $s_z>0$ and $C_z<\infty$ such that, for every
$0<s<s_z$ and almost every $0<h<h_z$,
\begin{equation}\label{eq:v65-seam-profile}
 \left|b_{z,s}(t_z+h)
       -\frac{J_z}{\pi}
          \arcsin\!\left(\min\left\{1,
                 \frac{s}{\kappa_z\sqrt{2h}}\right\}\right)\right|
            \le C_zJ_z h^{1/4}.
\end{equation}
The error constant does not depend on $s/h^{1/2}$, and so includes
both sides of the transition $h=s^2/(2\kappa_z^2)$.
For every fixed $s>0$ in this range,
\begin{equation}\label{eq:v65-half-jump}
 \lim_{h\downarrow0}b_{z,s}(t_z+h)=J_z/2
\end{equation}
for the representative defined by coarea. The restricted density
vanishes below $t_z$ within the chosen neighborhood.
The constants are local branch constants; no count-uniform bounds
for $h_z,s_z,C_z$ are asserted.
\end{theorem}
\begin{proof}
Translate $z$ to the origin in endpoint arclengths. The local action
is smooth with $dF(0)=0$ and positive Hessian $H$. A direct Morse
coordinate construction is useful here. Write
\[
 F(x)-t_z=\tfrac12x^{\mathsf T}Q(x)x,\qquad
 Q(x)=2\int_0^1(1-t)\nabla^2F(tx)\,dt.
\]
The positive square root of $Q(x)$ is smooth near zero. The map
$y=Q(x)^{1/2}x$ has derivative $H^{1/2}$ at zero, so it is a
local diffeomorphism with inverse $x=\Phi(y)$. It gives exactly
$F(\Phi(y))=t_z+|y|^2/2$. Choose a ball $|y|<r_z$ whose inverse
lies in $U$, and put $h_z=r_z^2/2$, reducing it when necessary.
This is the local neighborhood in the definition of $b_{z,s}$.

In these coordinates the original density is
$\widetilde w(y)=w(\Phi(y))|\det D\Phi(y)|$, with
$\widetilde w(0)=w_z/\sqrt{\det H}=J_z/(2\pi)$. Taylor's formula
gives, uniformly in $\theta$,
\[
 g(\Phi(r\omega_\theta))/r
       =a_z\cdot\omega_\theta+O_z(r),\qquad
 \widetilde w(r\omega_\theta)=J_z/(2\pi)+O_z(J_zr).
\]
The constants are relative to the fixed positive $J_z$; this places
no lower bound on $J_z$ uniformly in the count.
Polar coarea, including the Jacobian $r\,dr\,d\theta$, gives the
exact identity, with $r=\sqrt{2h}$,
\begin{equation}\label{eq:v65-polar-source}
 b_{z,s}(t_z+h)
    =\int_0^{2\pi}\widetilde w(r\omega_\theta)
             \one_{\{0<g(\Phi(r\omega_\theta))<s\}}\,d\theta.
\end{equation}
There is no additional factor $r^{-1}$: it cancels against the radial
Jacobian under $h=r^2/2$. Apply
Lemma~\ref{lem:v65-angular-threshold} with $\lambda=s/r$.
Its error is $O_z(r^{1/2})$, and the weight error is $O_z(J_zr)$.
Since $r^{1/2}=2^{1/4}h^{1/4}$, this proves
\eqref{eq:v65-seam-profile}. The angular formula is an actual density
representative. At fixed $s$, its leading angular measure tends to
$\pi$, proving \eqref{eq:v65-half-jump}. Positivity of the quadratic
normal form gives the assertion below $t_z$.
\end{proof}

\subsection{The unchanged smooth first-clearance source}
Let $q_0=47/53$ and $e'_j=q_0^{\min(j,m-1-j)/4}$. On the fixed
local neighborhood all the other physical margins are strictly
positive. For each fixed finite word, by reducing both the neighborhood
and $\varepsilon_z>0$, all their factors in $H_{m,R}^\varepsilon$
are therefore one when $0<\varepsilon<\varepsilon_z$. The sole
clearance factor equals zero for $g\le\varepsilon e'_j$ and one
for $g\ge2\varepsilon e'_j$, by the inherited guard convention.
The selected factor is consequently the first physical defect wherever
the local defect source is nonzero. No earlier factor or later incidence
is suppressed by a change of definition.

\begin{corollary}[Height of the original local defect]
\label{cor:v65-original-seam-source}
Put $s=\varepsilon e'_j$. Let $b^{\varepsilon,\mathrm{clr}}_z$ be
the density of the original first-clearance source restricted to the
local Morse disk above. For $0<\varepsilon<\varepsilon_z$ and
$0<h<h_z$, almost everywhere,
\begin{align}\label{eq:v65-guard-profile}
 \frac{J_z}{\pi}\arcsin\!\min\left\{1,
                       \frac{s}{\kappa_z\sqrt{2h}}\right\}
                   -C_zJ_zh^{1/4}
 &\le b^{\varepsilon,\mathrm{clr}}_z(t_z+h)\notag\\
 &\le\frac{J_z}{\pi}\arcsin\!\min\left\{1,
                       \frac{2s}{\kappa_z\sqrt{2h}}\right\}
                   +C_zJ_zh^{1/4}.
\end{align}
In particular the one-sided limit is $J_z/2$. On
$h\ge8s^2/\kappa_z^2$, the explicit leading upper term is at most
$C J_zs/(\kappa_z\sqrt h)$. For a bounded measurable insertion
$|w|\le M$, the variation density is bounded by $M$ times the same
positive local source.
\end{corollary}
\begin{proof}
On the original physical half-neighborhood,
\[
 \one_{\{0<g<s\}}\le1-H_{m,R}^\varepsilon
                         \le\one_{\{0<g<2s\}}.
\]
Apply positive pushforward and Theorem~\ref{thm:v65-seam-profile}
to the two bounds; this proves the inequalities and their common
one-sided limit. If $h\ge8s^2/\kappa_z^2$, the upper arcsine argument
is at most $1/2$, and $\arcsin x\le2x$ there. Variation domination
is performed on the original source before pushforward.
\end{proof}

\subsection{A positive weighted-cluster bound}
For fixed exact half-word labels $(n,k,m,R)$, let $\mathcal S$ be a
finite collection of the simple seam--minimum germs just described.
Use pairwise disjoint physical source neighborhoods; distinct word
copies are already disjoint. Each germ is counted once, with its
original responsible disk. Choose $0<r_z<h_z$ for each $z\in\mathcal S$.
Restrict its source further to $0<F-t_z<r_z$. For
$\varepsilon<\min_z\varepsilon_z$, let
$b^{\varepsilon,\mathcal S}$ be the sum of these restrictions.
The complete source then has the exact positive partition
\begin{equation}\label{eq:v65-positive-partition}
 b^{\varepsilon,\mathrm{clr}}
       =b^{\varepsilon,\mathcal S}
                         +b^{\varepsilon,\mathrm{outside}\ \mathcal S},
       \qquad b^{\varepsilon,\mathrm{outside}\ \mathcal S}\ge0.
\end{equation}
The second term includes every caustic point not assigned to these
neighborhoods, rather than only the transverse points.

\begin{theorem}[Critical-cluster control inside the selected caustic]
\label{thm:v65-cluster-height}
For almost every $t$,
\begin{align}\label{eq:v65-cluster-profile}
 b^{\varepsilon,\mathcal S}(t)
 &\le\sum_{z\in\mathcal S}J_z\one_{\{0<t-t_z<r_z\}}
 \left[\frac1\pi\arcsin\!\min\left\{1,
       \frac{2\varepsilon e'_{j_z}}{\kappa_z\sqrt{2(t-t_z)}}\right\}
                  +C_z(t-t_z)^{1/4}\right].
\end{align}
Consequently, if $r=\max_zr_z$ and
$\delta=\max_z C_zr_z^{1/4}$, the positive atomic measure of
\eqref{eq:v65-reversible-atomic-measure}, now formed from the specified
physical-side branch data, satisfies
\begin{equation}\label{eq:v65-cluster-height}
 b^{\varepsilon,\mathcal S}(t)
       \le\frac{1/2+\delta}{R c_*}
                           \mathfrak T_{n,k,m,R}([t-r,t]).
\end{equation}
These estimates retain coincident critical values and contain no
factor equal to the number of words or exact labels.
\end{theorem}
\begin{proof}
Sum the upper bound of Corollary~\ref{cor:v65-original-seam-source}
on the disjoint source neighborhoods. This proves
\eqref{eq:v65-cluster-profile}. The arcsine divided by $\pi$ is at
most $1/2$, and the other term is at most $\delta$. Every contributing
critical value lies in $[t-r,t]$. Substitute
\eqref{eq:v65-reversible-atomic-measure} to obtain
\eqref{eq:v65-cluster-height}. Addition of positive measures at an
equal value is exactly what the coarea sum requires.
\end{proof}

\begin{corollary}[An ordered transfer for a verified simple-seam family]
\label{cor:v65-ordered-trace-transfer}
For each fixed reconstruction band $B$, suppose a chosen family of
the preceding neighborhoods is valid for
$\varepsilon(B)=A_0B^{-1/12}$ at every sufficiently large count,
uniformly in the original exact labels and radius. Suppose its radii
satisfy $\sup_z r_z\le r_m\downarrow0$ and
$\sup_z C_zr_z^{1/4}\to0$, uniformly. If its positive reversible
measures satisfy
\begin{equation}\label{eq:v65-trace-concentration-input}
 \lim_{h\downarrow0}\limsup_{m\to\infty}
     \sup_{R,n,k,t}m^2\mathfrak T_{B,n,k,m,R}([t-h,t+h])=0,
\end{equation}
then $\sup_{R,n,k}m^2\|b^{\varepsilon(B),\mathcal S}\|_\infty\to0$
for that fixed $B$. The conclusion is also uniform over all
insertions of a fixed bounded modulus.
\end{corollary}
\begin{proof}
For fixed $h>0$, eventually $r_m<h$. Equation
\eqref{eq:v65-cluster-height} bounds the collision limsup by a fixed
multiple of that in \eqref{eq:v65-trace-concentration-input}.
Only after taking the collision limsup let $h\downarrow0$.
Positive source domination gives the insertion assertion. The band
has remained fixed throughout.
\end{proof}

The hypotheses of this last corollary are stated as inputs, not as
properties verified here for all Lorentz caustics. In particular, a
fixed $B$ does not automatically satisfy the local bound
$\varepsilon(B)<\varepsilon_z$ uniformly as $m$ increases. The
trace concentration and this neighborhood coverage both require
additional dynamical estimates. Multiple competing seams, vanishing
selected incidence, noncritical-roof tangencies, and section-boundary
junctions remain in the positive complement of
\eqref{eq:v65-positive-partition}. The first-incidence source is
unchanged. Substituting this partition into
Theorem~\ref{thm:v51-positive-error} preserves $P$, the signed
canonical arithmetic reference $G$, and both exact labels. It does
not yet remove that complement or establish the full pointwise law.

\paragraph{Proof route and comparison with the preceding estimates.}
The direct dependency route, within the compiled manuscript, is
\[
 \begin{gathered}
 \text{endpoint action and contact determinant (\S\S\ref{sec:v44-complete-height},
                                      \ref{sec:v61-inverse-incidence})}\\
 \Downarrow\\
 \text{reversible weight (Theorem~\ref{thm:v65-reversible-weight})}\\
 \Downarrow\\
 \text{physical-side Morse coarea (Theorem~\ref{thm:v65-seam-profile})}\\
 \Downarrow\\
 \text{exact-label positive cluster (Theorem~\ref{thm:v65-cluster-height}).}
 \end{gathered}
\]
The geometry of the image witness is supplied independently by
Lemma~\ref{lem:v49-physical-envelope}. Modules 131--136, including
their finite-type and buffered caustic estimates, remain intact.
The new profile controls a stated roof-critical germ within the tube,
where a pivot denominator is unavailable; it does not improve their
full-source count budgets by relabelling them. The arcsine profile
and the half coefficient are elementary consequences of quadratic
coarea. Their role here is to identify the actual physical coefficient
and its positive trace sum, rather than to claim a new general Morse
or Hill theorem.
